\begin{lemma}
For every $q\in\mathbb N$, the number of independent triples in the
neighborhood of $\ast$ in $L(A_q)$ is $3q^3$.
\end{lemma}
\begin{proof}
The definition \noderef{IndependentTripleCount} counts three-element sets
of neighbors with no adjacent distinct pair. By
\noderef{OrderThreeAffineGeometry} all these neighbors are nonvertical
indices, and by \noderef{LineGraphOfHypergraph} nonadjacency of distinct
indices means disjointness. The disjointness equivalence in
\noderef{OrderThreeAffineGeometry} forces every two indices to have equal
slopes and different intercepts. Consequently all three slopes agree,
and the three intercepts are exactly $0,1,2$.

Define the map
\[
(m,i_0,i_1,i_2)\longmapsto
\{(m,0,i_0),(m,1,i_1),(m,2,i_2)\}
\]
on $(\mathbb Z/3\mathbb Z)\times\{0,\ldots,q-1\}^3$.
Its three indices are distinct because their intercepts are distinct,
and their edges are pairwise disjoint by the geometry, so this is an
independent triple. For an arbitrary independent triple, its common slope
and the copy index of its unique element above each intercept $0,1,2$
give a unique inverse image. Reading off those data and then reconstructing
the set are inverse operations in both directions. This bijection proves
the exact count $3q^3$ for every $q$, with both sides empty when $q=0$.
\end{proof}
